\section{Storage, memory and distribution}
\label{sec:runtime}
The storage hierarchy makes the capacity strategy explicit
(Figure~\ref{fig:storage-hierarchy}): keep the complete relation and source
fields in a persistent store, stage one destination-row page through host
memory, and execute with a bounded temporary GPU working set. The final output
remains in device memory. FlashAttention's IO-aware view distinguishes on-chip
storage from device memory~\cite{flashattention}; here the paging boundary also
crosses host DRAM and NVMe. Tiga's paging mechanism bounds graph residency across
these outer memory levels; tile-local fusion bounds intermediate storage within a kernel.

\begin{figure}[htbp]
\centering
% Native vector diagram: widths are schematic, not measured capacities.
\begin{tikzpicture}[font=\small, >=Latex,
  tier/.style={draw=black!45, line width=.6pt, align=center, inner sep=6pt},
  note/.style={anchor=west, align=left, text width=4.3cm, font=\footnotesize},
  transfer/.style={->, line width=.8pt, draw=black!65}]
  \node[tier, minimum width=3.6cm, minimum height=1.0cm] (chip) at (4.5,6.5)
    {\textbf{On-chip storage}\\Registers / shared SRAM};
  \node[tier, draw=tigateal, fill=tigateal!7, line width=1.1pt,
    minimum width=6cm, minimum height=1.25cm] (gpu) at (4.5,4.4)
    {\textbf{GPU device memory}\\HBM or GDDR\\
     \textcolor{tigateal}{\textbf{Active page + full output}}};
  \node[tier, minimum width=7.5cm, minimum height=1.1cm] (ram) at (4.5,2.2)
    {\textbf{Host DRAM}\\File-backed pages + gather / staging buffers};
  \node[tier, minimum width=9cm, minimum height=1.1cm] (disk) at (4.5,0)
    {\textbf{NVMe-backed persistent store}\\Full CSR topology + source fields};

  \draw[transfer] (disk.north) -- node[right, font=\footnotesize]
    {Read / mmap; gather page inputs} (ram.south);
  \draw[transfer] (ram.north) -- node[right, font=\footnotesize]
    {Host-to-device page transfer} (gpu.south);
  \draw[<->, line width=.8pt, draw=black!65] (gpu.north) --
    node[right, font=\footnotesize] {Kernel loads / stores} (chip.south);

  \node[note] at (9.5,6.5) {Kernel-local working state\\No full edge-message array};
  \node[note, text=tigateal] at (9.5,4.4)
    {\textbf{Capacity constraint}\\Active page + full output\\must fit device memory};
  \node[note] at (9.5,2.2) {Bounded page staging\\Gather source values\\by destination-row page};
  \node[note] at (9.5,0) {Persistent graph storage\\Graph need not fit GPU\\Source fields stay on disk};
  \draw[black!30, densely dotted] (chip.east) -- (9.3,6.5);
  \draw[black!30, densely dotted] (gpu.east) -- (9.3,4.4);
  \draw[black!30, densely dotted] (ram.east) -- (9.3,2.2);
  \draw[black!30, densely dotted] (disk.east) -- (9.3,0);
\end{tikzpicture}

\caption{Storage hierarchy for paged graph execution. The central
constraint is \emph{active page plus full output}, not full-graph residency on
the GPU. Arrows show input movement and kernel access; the implementation also
returns each completed output page through host bytes before writing it into
the resident device output. No direct NVMe-to-GPU path is implied. Widths are
schematic, not hardware capacities or bandwidths. Experimental hardware and
measured footprints are reported in Section~\ref{sec:experimental-setup} and
Section~\ref{sec:billion}, respectively.}
\label{fig:storage-hierarchy}
\end{figure}

\subsection{Logical values versus physical residence}
\label{sec:gpu-paging}
The native runtime records value metadata and physical buffers separately.
Execution scopes select default device, tier budgets, spill directory and graph
paging defaults; they do not migrate preexisting values. Memory reports count
scope-created allocations, not process RSS or Torch-owned storage. Limits are
budgets, not reservations of physical capacity.

Ordinary Torch storage remains under Torch's allocator. Applications that need
Tiga-managed paging or distributed ownership use native \code{tg.Tensor}
fields for those paths, keeping the same message/reducer definition. The
\href{https://walkerchi.github.io/TIGA-lang/memory/}{memory API} and
\href{https://walkerchi.github.io/TIGA-lang/examples/distributed-memory/}{paged and distributed examples}
show the allocation scopes, graph snapshots and rank setup used by this runtime.

Temporary spill releases the tensor's resident-buffer reference while retaining
shape, dtype, logical device, version and same-process gradient history.
Observation restores the value. Other views or launches may still own storage,
so logical eviction is not proof that all associated device bytes were freed.
Persistent snapshots store values, not a reconstructible autograd graph. Device
copy is differentiable on supported native paths but currently eager, not merely
a scheduler hint. The optional LRU policy evicts whole idle native tensors; it
does not tile a working set larger than device memory, manage Torch allocations
or supply transparent GPU-to-NVMe execution for arbitrary kernels.

Paged CSR streams bounded destination-row ranges on CPU and, for the tested
FP32 forward slice, on CUDA via host staging. CPU paging has a VJP path;
CUDA paged backward remains unsupported. Graph offload and
field storage have distinct memory lifetimes; automatically paging topology
does not imply every node/edge field is paged. Forward output assembly allocates
the final resident tensor once and writes each completed page into its row
range. Only one output-page payload is retained on the host; the final output
must still fit the selected device. Persistent CSR ingestion in the large-graph
experiment writes bounded chunks directly to the versioned store, without
constructing a graph-sized Python index list.

This bounded-memory path is currently a capacity mechanism rather than an
optimal IO pipeline. Source fields are gathered on the host into page-edge
order and copied to CUDA, so repeated source IDs can duplicate feature traffic
within a page. Each completed output page is copied back to host bytes before
being written into the final device output. Page-local staging bounds the
temporary working set but does not eliminate these copies. Section~\ref{sec:billion} measures their
cost and separates it from GPU kernel execution. Deduplicating source staging,
direct device output assembly and asynchronous IO are distinct optimizations;
they are not implied by an offload interface.

The optional runtime prefetcher maintains a bounded queue of topology-page
reads and contiguous-field read-ahead hints. It does not move source gathers,
host packing or CUDA copies into a compiler-scheduled asynchronous pipeline.
Consequently prefetch depth is not a guarantee of disk/PCIe/kernel overlap;
the capacity experiment reports the explicit serialized configuration.

\subsection{Ownership, halos and progress}
A rank is a cooperating worker process. A logical device mesh and destination
partition assign output ownership to ranks. Remote source values needed by local
destinations are held as ghost copies; the halo map specifies which values must
cross rank boundaries. The public schedule completes halo exchange before computing all owned
destination rows in one local call, without split-output assembly.
Reverse-mode execution accumulates ghost contributions back
to owners. Event ordering is a correctness requirement independently of whether
communication and computation actually overlap in time.
For example, if a destination on rank 0 reads a source owned by rank 1, the
forward exchange copies that source value to rank 0. Backward sends the source's
gradient contribution in the opposite direction and adds it to the gradient on
rank 1. This is why ownership must be retained through differentiation, not
treated solely as a forward communication detail.

The runtime supports CPU rank-local execution, MPI transport, and two-host CUDA
forward/VJP through TCP and NCCL. Deployment supplies worker processes, rank
membership and a communicator. Ownership remains fixed during execution;
Section~\ref{sec:two-host} evaluates an equal-destination partition on spatial meshes.

\subsection{Validation and trust boundaries}
Dispatch distinguishes supported compiled paths from reference/provider routes.
Full CSR validation checks monotonic row pointers, endpoints and source-index
bounds; basic validation checks metadata and is not an exhaustive scan.
Paged reads validate each consumed range. Tensor snapshots, geometry importers
and TCP transports are not security sandboxes: compiled UDFs and providers
execute trusted native code, and transport deployment requires trusted peers.
\FloatBarrier
